\documentclass[12pt,a4paper]{article}
\usepackage[utf8]{inputenc}
\usepackage[T1]{fontenc}
\usepackage[brazil]{babel}
\usepackage{indentfirst}
\usepackage[margin=2.5cm]{geometry}

\title{Resenha Cr\'itica: Artigo 12 ---\\
Gera\c{c}\~ao de C\'odigo DSL para Diagramas C4 Model via LLMs}
\author{Gustavo Pessoa Firmino Duarte}
\date{23 de Abril de 2026}

\begin{document}
\maketitle

\section{Introdu\c{c}\~ao}

O artigo ``Gera\c{c}\~ao de C\'odigo DSL para Diagramas C4 Model via LLMs''
explora a interseção entre dois avan\c{c}os recentes e complementares: o modelo
C4 para visualiza\c{c}\~ao de arquitetura de software e os \textit{Large Language
Models} (LLMs) para gera\c{c}\~ao autom\'atica de c\'odigo. O trabalho investiga
a capacidade de modelos de linguagem como o GPT-4 de gerar, a partir de descri\c{c}\~oes
em linguagem natural, c\'odigo em DSLs (\textit{Domain-Specific Languages}) como
o Structurizr DSL ou o PlantUML, que produzem diagramas C4 v\'alidos e semanticamente
coerentes. A proposta \'e particularmente relevante porque documenta\c{c}\~ao
arquitetural \'e frequentemente negligenciada por seu custo em tempo, e LLMs
poderiam reduzir substancialmente essa barreira.

\section{Abordagem, Resultados e Limita\c{c}\~oes}

A metodologia central do artigo consiste em avaliar se LLMs conseguem transformar
descri\c{c}\~oes textuais de sistemas de software em c\'odigo DSL que produza
diagramas C4 sintaticamente corretos e semanticamente representativos. Os
experimentos consideram diferentes n\'iveis da hierarquia C4 (especialmente
Context e Container), variando a especificidade das descri\c{c}\~oes fornecidas
ao modelo e avaliando a qualidade dos resultados por crit\'erios como
\textit{validade sint\'atica}, \textit{completude dos elementos} e \textit{fidelidade
sem\^antica} em rela\c{c}\~ao \`a descri\c{c}\~ao original.

Os resultados indicam que LLMs s\~ao capazes de gerar c\'odigo DSL sintaticamente
correto para casos relativamente simples, especialmente quando a descri\c{c}\~ao
do sistema \'e detalhada e estruturada. Por\'em, as limita\c{c}\~oes s\~ao
significativas: descri\c{c}\~oes amb\'iguas ou sistemas complexos com muitas
depend\^encias tendem a produzir diagramas que omitem rela\c{c}\~oes importantes
ou introduzem elementos n\~ao mencionados. Al\'em disso, a fidelidade sem\^antica
---  garantir que o diagrama representa fielmente a arquitetura real e n\~ao apenas
uma interpreta\c{c}\~ao plaus\'ivel --- permanece um desafio aberto.

O artigo tamb\'em discute \textit{prompt engineering} como fator cr\'itico:
prompts bem estruturados, que explicitam a hierarquia C4 esperada e o vocabul\'ario
da DSL alvo, produzem resultados substancialmente melhores. Isso sugere que a
integra\c{c}\~ao de LLMs em fluxos de documenta\c{c}\~ao arquitetural n\~ao
elimina o trabalho humano, mas o transforma: em vez de escrever DSL manualmente,
o arquiteto passa a formular descri\c{c}\~oes precisas e a revisar e corrigir o
c\'odigo gerado.

\section{Conex\~ao com a Pr\'atica Profissional}

Este artigo descreve uma pr\'atica que j\'a adoto informalmente: o uso de LLMs
para acelerar tarefas de desenvolvimento que envolvem geração de c\'odigo a partir
de especifica\c{c}\~oes textuais. No projeto acad\^emico de cinco dias com
tecnologias desconhecidas, usei ferramentas de IA extensivamente --- n\~ao para
gerar diagramas, mas para scaffolding de c\'odigo, compreensi\~ao de APIs e
depura\c{c}\~ao. A experi\^encia confirmou a observa\c{c}\~ao central do artigo:
o valor da IA est\'a na acelera\c{c}\~ao, mas a qualidade do resultado depende
criticamente da qualidade da especifica\c{c}\~ao fornecida pelo humano. Prompts
vagos produzem c\'odigo gen\'erico; prompts precisos produzem c\'odigo acion\'avel.

A ideia de usar LLMs para gerar diagramas C4 \'e algo que pretendo experimentar
nos pr\'oximos projetos: descrever o sistema em linguagem natural, gerar o
Structurizr DSL e refinar iterativamente. O ciclo de feedback \'e mais r\'apido
do que escrever DSL manualmente e o resultado \'e um artefato de documenta\c{c}\~ao
versionado e mant\'ivel.

\section{Conclus\~ao}

O artigo faz uma contribui\c{c}\~ao pr\'atica e oportuna ao investigar a
aplicabilidade de LLMs a uma tarefa espec\'ifica e de alto valor: a gera\c{c}\~ao
autom\'atica de documenta\c{c}\~ao arquitetural em formato C4. Os resultados s\~ao
promissores para casos simples e apontam claramente as condi\c{c}\~oes necess\'arias
para resultados confi\'aveis: descri\c{c}\~oes precisas, prompts estruturados e
revis\~ao humana. A principal lic\~ao \'e que LLMs n\~ao substituem o conhecimento
arquitetural, mas podem democratizar a documenta\c{c}\~ao --- tornando mais
acess\'ivel a equipes menores sem arquitetos dedicados a produ\c{c}\~ao de
diagramas claros e atualizados.

\end{document}
